\documentclass[10pt,a4paper]{amsart}
\usepackage[margin=19mm]{geometry}
\usepackage{amsmath,amssymb,mathtools}
\usepackage[scr=rsfs,cal=euler]{mathalfa}
\usepackage{xfrac}
\usepackage[unicode,hidelinks,bookmarks=false]{hyperref}
\newcommand{\ie}{i.e.\ }
\newcommand{\cf}{cf.\ }
\newcommand{\resp}{respectively\ }
\newcommand{\Subsection}{\subsection}
\providecommand{\nobreakdash}{\nobreak\hbox{-}}
\setlength{\emergencystretch}{2em}
\multlinegap0pt
\usepackage{amssymb}
\usepackage{enumerate}
\usepackage{xcolor}
\usepackage{mathtools}
\newtheorem{proposition}{Proposition}
\title{Unified growth rates for operator semigroups under generalized Kreiss conditions}
\author{Filippo Dell'Oro, Alberto Farina, Vittorino Pata}
\date{Source: arXiv:2608.16217v1 (CC BY 4.0)}
\begin{document}
\maketitle

\noindent\emph{Source and licence.} Unified growth rates for operator semigroups under generalized Kreiss conditions. Version arXiv:2608.16217v1 (CC BY 4.0), distributed by arXiv under the Creative Commons Attribution 4.0 International licence, http://creativecommons.org/licenses/by/4.0/, as recorded in the arXiv licence metadata for this exact version and shown on the abstract page https://arxiv.org/abs/2608.16217v1. Full licence notice: \texttt{licenses/arxiv-2608.16217-CC-BY-4.0.txt}.\par\smallskip

\noindent\textit{Editorial scope.} Counted unit: the article's proposition coro (source0.tex lines 205--220). The article's complete original proof of it is delivered with it (source0.tex lines 222--229, one \texttt{proof} environment of 8 lines). No mathematics is written, repaired or re-derived: the statement and the proof are the article's own lines, in the article's own notation, and no retained line is substituted or re-flowed.  No further source-local context is needed: every cross-reference of every retained line resolves inside the two delivered spans.  Nothing was omitted from any retained span, so the package carries no omission record.  No retained line contains a citation, so the wrapper carries no bibliography and no bibliography entry.  No retained line contains a figure environment, an image-inclusion command or drawing code, so no image is delivered and the package declares no asset dependency.  Editorial formatting only, in addition: an AMS \texttt{amsart} preamble that restates the article's own declarations and macros used by the retained lines, namely the environments \texttt{proposition} on separate counters, together with the standard packages those lines need.  The retained environments print in this document as Proposition 1 in order of appearance, whereas the article prints the same items with the numbering of its own full document.  The complete native source of the article as distributed in the arXiv e-print for this exact version is bundled unchanged as \texttt{sources/207813/source0.tex}.
\begin{proposition}
\label{coro}
The following hold:
\begin{enumerate}
\item[\rm (i)] We always have
$$\frac{g(t)}{\sqrt{\Lambda(t)}}=O(g(t)).$$
In particular,
$$\frac{g(t)}{\sqrt{\Lambda(t)}}=o(g(t))$$
whenever $\Lambda(t)\to\infty$, i.e., whenever $t/[g(t)]^2$ is not integrable at infinity.
\vskip2mm
\item[\rm (ii)] If $h$ is non-decreasing we have
$$
\frac{g(t)}{\sqrt{\Lambda(t)}} = O\left(g(t)\frac{h(t)}{\sqrt{\log t}}\right).
$$
\end{enumerate}
\end{proposition}

\begin{proof}
As $\Lambda$ is increasing, point (i) is immediate.
To show point (ii), we exploit the monotonicity of $h$.
Then, for $2<s<t$, we write $g(s)\leq sh(t)$. Hence,
$$\Lambda(t)\geq \frac{1}{[h(t)]^2} \int_2^t\frac{1}{s}\,ds
\sim \frac{\log t}{[h(t)]^2},$$
yielding the desired conclusion.
\end{proof}

\end{document}
